% ============================================================
%  paper-export :: 公共 preamble（gb / report 两种预设共用）
%
%  职责：字体、符号兜底、表格、列表、代码块、引用块、分页质量。
%  版式相关（页面尺寸、标题层级、页眉页脚、目录）在 preamble-<style>.tex。
%
%  注入点：pandoc --include-in-header，位于所有 \usepackage 之后，
%          但在 hyperref 之前 —— 因此这里不能调用 \hypersetup。
% ============================================================

% ---------- 0. 全文黑色 ----------
% 硬性要求：除图片与极浅的底色外，一切文字纯黑。
% 链接颜色由 export.sh 通过 -V linkcolor=black 等传入（hyperref 加载晚于此处）。
\usepackage{xcolor}
\definecolor{pblack}{HTML}{000000}
\definecolor{prule}{HTML}{000000}       % 表格线与页眉线
\definecolor{pcodebg}{HTML}{F7F7F8}     % 代码块底色（唯一允许的非黑元素之一）
\definecolor{pcodeframe}{HTML}{E3E3E6}
\definecolor{pquotebar}{HTML}{BFBFBF}
\definecolor{ptablehead}{HTML}{F2F2F2}
\definecolor{ptablesoft}{HTML}{D9D9D9}
\AtBeginDocument{\color{pblack}}

% ---------- 1. 中文字体族 ----------
% ctex 在 macOS 上默认把 \heiti 落到 STXihei（华文细黑）：笔画细、字面小，
% 三号标题看上去反而比小四正文还弱，层级感全无。显式换成中等字重的黑体。
% \kaishu 供封面信息栏使用。
\setCJKfamilyfont{zhhei}{Heiti SC Medium}[BoldFont=Heiti SC Medium]
\newfontfamily\peheiti{Heiti SC Medium}[BoldFont=Heiti SC Medium]
\renewcommand{\heiti}{\peheiti\CJKfamily{zhhei}}
\setCJKfamilyfont{zhsong}{Songti SC}
\newfontfamily\pesongti{Songti SC}
\renewcommand{\songti}{\pesongti\CJKfamily{zhsong}}
\setmainfont{Songti SC}
% 字体命令同时切换中西文，避免标题或封面里的英文突然落到另一套西文字体。
\IfFontExistsTF{Kaiti SC}{%
  \setCJKfamilyfont{zhkai}{Kaiti SC}
  \newfontfamily\pekaishu{Kaiti SC}
  \renewcommand{\kaishu}{\pekaishu\CJKfamily{zhkai}}
}{%
  \renewcommand{\kaishu}{\pesongti\CJKfamily{zhsong}}
}

% 中西文之间自动留出细气隙：中文技术文档里「AI 拆解」「H5 前台」这类混排极多，
% 不留气隙会粘连，留全角又太散。
\xeCJKsetup{xCJKecglue={\hskip 0.08em plus 0.03em minus 0.02em}}

% ---------- 2. 符号字形兜底 ----------
% 正文字体（宋体 / Times）不含 Box Drawing 与几何符号，缺字形会静默丢字。
% Menlo 出自 DejaVu 血统，U+2500–257F、U+25A0–25FF、U+2600–26FF 覆盖完整，
% 用它兜底这些字符，视觉上与正文差异极小。
\usepackage{newunicodechar}
% 缩放必须与 \setmonofont 完全一致，否则这些符号在代码块里与相邻等宽字符的
% 字身宽对不上，横线会出现缝隙、框线接不拢。
%
% Maple Mono NF CN 是 Nerd Font 版，Box Drawing / 几何符号 / 箭头覆盖比 Menlo
% 更全，且与正文等宽字体同源，混排时字重和字面一致。没装则回退 Menlo。
\IfFontExistsTF{Maple Mono NF CN}
  {\newfontfamily\symbolfont{Maple Mono NF CN}[Scale=0.90]}
  {\newfontfamily\symbolfont{Menlo}[Scale=0.85]}
% 上面那句「覆盖比 Menlo 更全」有两个例外：Maple Mono NF CN Regular 里
% ★ U+2605 与 ☆ U+2606 **没有字形**，落到 symbolfont 就静默丢字
% （tectonic 只报一行 Missing character，PDF 里那个位置是空的）。
% 把这两个字符单独指到 Menlo —— 逐个编译验证过，其余 43 个符号 Maple 都有。
\IfFontExistsTF{Menlo}
  {\newfontfamily\symbolaltfont{Menlo}[Scale=0.85]}
  {\let\symbolaltfont\symbolfont}
\makeatletter
% 替换体必须用码点 \char"XXXX 而非字符本身：newunicodechar 定义的是
% active character，在自己的替换体里再写一遍该字符会自指，导致字形丢失。
\def\pe@sym#1#2{\newunicodechar{#1}{{\symbolfont\char"#2\relax}}}
% Maple 缺字形的那几个走备用字体
\def\pe@symalt#1#2{\newunicodechar{#1}{{\symbolaltfont\char"#2\relax}}}
% 制表符（ASCII 框图）
\pe@sym{─}{2500}\pe@sym{━}{2501}\pe@sym{│}{2502}\pe@sym{┃}{2503}
\pe@sym{┌}{250C}\pe@sym{┐}{2510}\pe@sym{└}{2514}\pe@sym{┘}{2518}
\pe@sym{├}{251C}\pe@sym{┤}{2524}\pe@sym{┬}{252C}\pe@sym{┴}{2534}\pe@sym{┼}{253C}
\pe@sym{═}{2550}\pe@sym{║}{2551}\pe@sym{╔}{2554}\pe@sym{╗}{2557}
\pe@sym{╚}{255A}\pe@sym{╝}{255D}
% 几何与标记符号
\pe@sym{■}{25A0}\pe@sym{□}{25A1}\pe@sym{●}{25CF}\pe@sym{○}{25CB}
\pe@sym{◆}{25C6}\pe@sym{◇}{25C7}\pe@sym{▲}{25B2}\pe@sym{▼}{25BC}
\pe@sym{►}{25BA}\pe@sym{◄}{25C4}\pe@symalt{★}{2605}\pe@symalt{☆}{2606}
\pe@sym{☐}{2610}\pe@sym{☑}{2611}\pe@sym{✓}{2713}\pe@sym{✔}{2714}
\pe@sym{✗}{2717}\pe@sym{✘}{2718}\pe@sym{⚠}{26A0}
% 箭头
\pe@sym{→}{2192}\pe@sym{←}{2190}\pe@sym{↑}{2191}\pe@sym{↓}{2193}
\pe@sym{⇒}{21D2}\pe@sym{⇐}{21D0}\pe@sym{↔}{2194}
\makeatother

% 带圈数字 ①②③，供三级有序列表使用。
% 注意 Songti SC **没有** U+2460 区段（实测 fc-list :charset=2460 里没有它），
% 直接 \songti\char"2460 会排出一个空方框；Heiti SC 才有。
% 用 \newfontfamily 显式绑定，绕开 xeCJK 的字符分类。
\newfontfamily\pecirclefont{Heiti SC}
\newcommand{\pecircled}[1]{{\pecirclefont\char\numexpr"2460+#1-1\relax}}

% 制表符（Box Drawing / Block Elements）声明为非 CJK 字符。
% 上面那套 \newunicodechar 是 active character，而 verbatim 环境里所有字符都是
% catcode 12，active 不展开 —— 代码块里的框线根本走不到兜底，会落到
% \setCJKmonofont（宋体，全角宽）上，于是 1 格宽的横线被排进 2 格，
% 框线断成一截一截。声明成 Default 类后它们改走 Menlo，宽度恰好 1 格。
\xeCJKDeclareCharClass{Default}{"2500 -> "259F}

% ---------- 3. 代码块 / ASCII 框图 ----------
% 拉丁等宽用 Menlo；CJK 等宽必须恰好是拉丁等宽的 2 倍宽，
% 终端里画的框图（中文占 2 格）才能在 PDF 中保持对齐。
% Maple Mono CN 一套字体同时供拉丁与 CJK，中文字身宽恰好是拉丁的 2 倍 ——
% 这正是终端画 ASCII 框图的前提，比「Menlo 0.85 + 宋体 1.02」这种靠缩放
% 凑出来的 2:1 可靠得多（后者的比例只在特定字号下准确）。
% 没装则退回原来的拼接方案。
\IfFontExistsTF{Maple Mono CN}
  {\setmonofont{Maple Mono CN}[Scale=0.90]
   \setCJKmonofont{Maple Mono CN}[Scale=0.90]}
  {\setmonofont{Menlo}[Scale=0.85]
   \setCJKmonofont{Songti SC}[Scale=1.02]}

% ---------- 3b. 行内代码里的 CJK 跟随正文 ----------
% 上面的 \setCJKmonofont 同时作用于代码块与行内代码（\texttt）。对代码块这是
% 对的 —— 中文必须与拉丁保持 2:1 才能让 ASCII 框图对齐。但行内代码是嵌在正文
% 句子里的：一句「来源 `docs/旧版pacs 模块/接口规范.docx` 经交叉验证」会在同一
% 行里并排出现两种中文字体（等宽 vs 正文），字形与字距都对不上。技术文档里带
% 中文的路径、表名、菜单名很常见，整页看下来就是「中英文字体不一致」的观感。
%
% 只把 \texttt 内的 CJK 切回正文族：拉丁仍走等宽（保住 pacs_exam_order 这类
% 标识符的辨识度与下划线可读性），中文与正文一致。代码块走 \ttfamily /
% Verbatim，不经过 \texttt，因此完全不受影响，框图的 2:1 对齐照旧。
%
% 用 \DeclareTextFontCommand 而不是 \renewcommand：前者是 LaTeX 声明文本字体
% 命令的正规入口，保留 \texttt 的 robust 性与 \protect 行为。这里只在内容之前
% 切字体族，不碰内容 token —— 与 §8 注释里「往内容里插 \penalty 会喂坏 pandoc
% 生成的转义序列」那个坑不是一回事，那条禁令依然有效。
%
% \CJKrmdefault 是 xeCJK 为「CJK 正文族」保留的族名，\setCJKmainfont 写的就是
% 它。因此这里不硬编码任何字体名，modern 覆盖成普惠体、report 保持宋体、
% brief 用仿宋，各预设都会自动跟随各自的正文字体。
\DeclareTextFontCommand{\texttt}{\ttfamily\CJKfamily{\CJKrmdefault}}

\usepackage{fvextra}
\fvset{breaklines=true,breakanywhere=true,fontsize=\zihao{5}}

% pandoc 对「无语言标注」的代码块（裸围栏，技术文档里的 ASCII 框图几乎都是
% 这一类）生成原生 verbatim，既不受 \fvset 控制也不会断行，超宽必然溢出版心。
% 用 fancyvrb 接管：统一五号字，且只在空格处断行——框图不会被从中间劈开。
\RecustomVerbatimEnvironment{verbatim}{Verbatim}%
  {fontsize=\zihao{5},breaklines=true,breakanywhere=false,xleftmargin=0pt}

% 带语言标注的代码块走 pandoc 的 Highlighting 环境，它用 commandchars 把每个词
% 包进 \NormalTok{} 之类的宏里。fvextra 默认不在宏参数内部断行，一条无空格的
% 长 URL 会整条溢出版心。重新定义并显式打开 breakanywhere。
% 必须在 \AtBeginDocument 里做 —— pandoc 的定义写在 preamble 更靠后的位置。
\usepackage{mdframed}
% 注意 \makeatletter：\AtBeginDocument 在「读到」时就把 token 存下来，
% 里面的 \@ifundefined 若在 @ 为普通字符时读入，会被拆成 \@ + ifundefined，
% 在 \begin{document} 处报「You can't use \spacefactor in vertical mode」。
\makeatletter
\AtBeginDocument{%
  \DefineVerbatimEnvironment{Highlighting}{Verbatim}%
    {commandchars=\\\{\},breaklines=true,breakanywhere=true,fontsize=\zihao{5}}%
  % 代码块外框：极浅灰底 + 细边框，取代 pandoc 默认的无框 snugshade
  \@ifundefined{Shaded}{}{%
    \renewenvironment{Shaded}%
      {\begin{mdframed}[backgroundcolor=pcodebg,linecolor=pcodeframe,
                        linewidth=0.4pt,roundcorner=2pt,
                        innerleftmargin=8pt,innerrightmargin=8pt,
                        innertopmargin=5pt,innerbottommargin=5pt,
                        skipabove=7pt,skipbelow=7pt]}%
      {\end{mdframed}}}%
}
\makeatother
% 代码块内不继承正文的 1.5 倍行距
\AtBeginEnvironment{Verbatim}{\linespread{1.0}\selectfont}

% ---------- 4. 表格 ----------
% 规范：三线表，无竖线；表格居中；五号字；行距 1.0；跨页重复表头。
% 列宽由 assets/tablefit.lua 按内容重算后写进 AST，此处只管观感。
\usepackage{etoolbox}
\usepackage{array}
\usepackage{booktabs}
\usepackage{longtable}
\usepackage{ragged2e}
% colortbl 必须在 \arrayrulecolor 之前加载
\usepackage{colortbl}

% 三线表线宽：顶/底线与表头下线是「结构线」，纯黑；
% 行间分界线是「阅读辅助线」，必须细而浅 —— 和结构线一样粗一样黑的话，
% 整张表会退化成网格，比不画还吵。
\setlength{\heavyrulewidth}{0.9pt}
\setlength{\lightrulewidth}{0.4pt}
\setlength{\cmidrulewidth}{0.3pt}
\arrayrulecolor{prule}

% 行间分界线：由 tablefit.lua 在「有单元格折行」的表里逐行插入
% \noalign{\perowrule}。线宽与颜色集中在这里调。
\definecolor{prowrule}{HTML}{D0D0D0}
\newcommand{\perowrule}{{\color{prowrule}\hrule height 0.3pt}}

% 单元格上下留白：靠 \arraystretch 而非行距，避免多行单元格被撑得过高
\renewcommand{\arraystretch}{1.18}
\setlength{\tabcolsep}{5pt}

% longtable 水平居中（pandoc 生成的 @{} 会让它默认贴左）
\setlength{\LTleft}{0pt plus 1fill}
\setlength{\LTright}{0pt plus 1fill}
\setlength{\LTpre}{8pt}
\setlength{\LTpost}{8pt}

% 表格整体五号字、行距 1.0：正文的 1.5 倍行距套进表格会让每一行都过高
\AtBeginEnvironment{longtable}{\zihao{5}\linespread{1.05}\selectfont}
\AtBeginEnvironment{tabular}{\zihao{5}\linespread{1.05}\selectfont}

% ---------- 5. 列表 ----------
% 规范：无序 • / – / · ；有序 1. / (1) / ① ；缩进与正文首行对齐；项间紧凑。
\usepackage{enumitem}
% 一级列表的标记与正文首行缩进（2em）对齐 —— 标记顶在版心左边缘会让列表
% 看上去比正文还往外凸。二三级各再进 2em。
% labelwidth 显式给出，不靠 enumitem 从 leftmargin 反推 —— 反推值会把标记
% 与文字之间撑出一大段空白。三个量的关系：
%   labelindent(标记左边缘) + labelwidth + labelsep = leftmargin(文字左边缘)
\setlist{topsep=3pt,partopsep=0pt,parsep=0pt,itemsep=2pt,align=left}
\setlist[itemize,1]{label=\textbullet,
                    labelindent=2em,labelwidth=0.7em,labelsep=0.4em,leftmargin=3.1em}
\setlist[itemize,2]{label=\textendash,
                    labelindent=0em,labelwidth=0.7em,labelsep=0.4em,leftmargin=1.6em}
\setlist[itemize,3]{label=\textperiodcentered,
                    labelindent=0em,labelwidth=0.7em,labelsep=0.4em,leftmargin=1.6em}
\setlist[enumerate,1]{label=\arabic*.,
                    labelindent=2em,labelwidth=1.1em,labelsep=0.4em,leftmargin=3.5em}
\setlist[enumerate,2]{label=(\arabic*),
                    labelindent=0em,labelwidth=1.6em,labelsep=0.4em,leftmargin=2.2em}
\setlist[enumerate,3]{label=\pecircled{\value{enumiii}},
                    labelindent=0em,labelwidth=1.1em,labelsep=0.4em,leftmargin=1.8em}
% pandoc 的 tightlist 会把项间距清零；这里保留一点点，避免中文行距下挤成一坨
\providecommand{\tightlist}{\setlength{\itemsep}{1pt}\setlength{\parskip}{0pt}}

% 列表内不再首行缩进（正文缩进 2em 已由 leftmargin 承担）
\AtBeginEnvironment{itemize}{\setlength{\parindent}{0pt}}
\AtBeginEnvironment{enumerate}{\setlength{\parindent}{0pt}}
\AtBeginEnvironment{description}{\setlength{\parindent}{0pt}}

% ---------- 6. 引用块 ----------
% Markdown 的 > 引用：左侧竖线 + 极浅灰底
\renewenvironment{quote}{%
  \begin{mdframed}[linewidth=3pt,linecolor=pquotebar,
                   topline=false,bottomline=false,rightline=false,
                   backgroundcolor=pcodebg,
                   innerleftmargin=10pt,innerrightmargin=10pt,
                   innertopmargin=6pt,innerbottommargin=6pt,
                   skipabove=6pt,skipbelow=6pt]\setlength{\parindent}{0pt}}%
  {\end{mdframed}}

% ---------- 7. 图片 ----------
% 超出版心的图自动缩到版心宽；居中；不产生「Overfull \hbox」
\usepackage{graphicx}
\setkeys{Gin}{keepaspectratio}
\makeatletter
\def\maxwidth{\ifdim\Gin@nat@width>\linewidth\linewidth\else\Gin@nat@width\fi}
\makeatother

% ---------- 8. 断词与溢出 ----------
% 注：行内代码（\texttt）的断行由 sanitize.lua 在 AST 层处理，不在这里
% 重定义 \texttt —— 那样会把 \penalty 喂给 pandoc 生成的转义序列
% （如 ^ → \^{}），导致 "Missing number" 编译中断。
\sloppy
\emergencystretch=3em
\hbadness=10000  % 抑制 underfull hbox 噪音警告

% ---------- 9. 分页质量 ----------
% 孤行/寡行：段落首行不许单独留在页底，末行不许单独跑到下页
\widowpenalty=10000
\clubpenalty=10000
\displaywidowpenalty=10000
\brokenpenalty=10000
% 标题后至少要能容下 3 行正文，否则整个标题推到下一页 ——
% 杜绝「标题孤零零挂在页面最后一行」
\usepackage{needspace}
\pretocmd{\section}{\needspace{4\baselineskip}}{}{}
\pretocmd{\subsection}{\needspace{4.5\baselineskip}}{}{}
\pretocmd{\subsubsection}{\needspace{4\baselineskip}}{}{}
\pretocmd{\paragraph}{\needspace{3.5\baselineskip}}{}{}
% 标题后紧跟表格是中文技术文档的常态。LaTeX 只知道「标题后有内容」，
% 于是「（六）术语说明」留在页尾、表格整个甩到下一页。给 longtable 也要一段
% 起始空间，标题与表头才会待在同一页。
\BeforeBeginEnvironment{longtable}{\needspace{4\baselineskip}}
\usepackage{float}

% ---------- 10. 目录观感 ----------
% 默认点线过疏（\@dotsep=4.5mu），中文目录看上去像断续的省略号；收紧到 1.8。
% 各级条目的字体与缩进在 preamble-<style>.tex 里按预设分别设定。
\makeatletter
\renewcommand{\@dotsep}{1.8}
\makeatother
